Table~\ref{tab:ablation} removes or changes one component of E-DDQN at a time while keeping the configuration selected for each outer fold fixed (five seeds; same protocol). No formal test was performed, so the differences below are point estimates. Because the selected configurations already used no warm-up in four of the five folds, the ``no warm-up'' variant differs from the full model only in the remaining fold. This analysis and the reward-ratio analysis below hold E-DDQN's primary (draw-1) configurations fixed, which rank weakly (Section~\ref{sec:why}); they cannot show whether a component helps a configuration that ranks well. Every variant ranks seconds only weakly (ROC-AUC \AblSoftCostRocAuc--\AblGammaZeroCostRocAuc, where 0.5 is chance). No single change moves the cost by more than about one seed standard deviation of the full model (\EDDQNCostCostMean$\pm$\EDDQNCostCostSD). The lowest cost is obtained by \emph{replacing} the asymmetric reward with the symmetric $\pm1$ reward (\AblAsymCostCost) and the highest ROC-AUC by switching off bootstrapping ($\gamma=0$; \AblGammaZeroCostRocAuc); hard target copies (\AblSoftCostCost), no dueling head (\AblDuelCostCost), uniform replay (\AblBalCostCost), no warm-up (\AblWarmCostCost) and toggling full-information replay (\AblFullCostCost) leave the cost essentially unchanged, and removing the double estimator slightly increases it (\AblDoubleCostCost). None of these components can be credited with a gain in safety or accuracy. Every variant is scored on the same nominal implied-probability scale and thresholded with the same rules, so the reward table can act only through the learned scores; at the fixed cut-off $\hat p\ge1/6$ (rule D0) the symmetric-reward variant is not worse either (cost \AblAsymBayesCost{} vs.\ \EDDQNBayesCostMean), although the effective position of that cut-off on $P(y=1\mid s)$ differs between the variants (Section~\ref{sec:formulation}).

\begin{table*}[t]
\centering
\caption{Leave-one-component-out ablation of E-DDQN (selected configuration per outer fold held fixed; rule D1; mean over 5 seeds, cost $\pm$ SD). The no-warm-up variant changes only the one fold whose selected configuration used warm-up.}
\label{tab:ablation}
\footnotesize
\begin{tabular}{lccccc}
\toprule
Variant & Recall & Spec. & Cost$_{5:1}$ & ROC-AUC & PR-AUC \\
\midrule
Full E-DDQN (selected) & 93.4 & 7.4 & 0.663$\pm$0.048 & 0.538 & 0.483 \\
\midrule
full-information replay toggled & 93.7 & 5.4 & 0.667$\pm$0.106 & 0.563 & 0.500 \\
symmetric $\pm$1 reward instead of asymmetric & 94.8 & 8.6 & 0.625$\pm$0.038 & 0.572 & 0.517 \\
uniform instead of class-balanced replay & 94.3 & 4.3 & 0.660$\pm$0.045 & 0.554 & 0.486 \\
no double estimator (max-operator target) & 93.1 & 5.3 & 0.681$\pm$0.086 & 0.543 & 0.491 \\
no dueling head & 94.6 & 5.8 & 0.644$\pm$0.059 & 0.559 & 0.494 \\
hard target copy every 10 updates & 94.9 & 6.5 & 0.634$\pm$0.043 & 0.529 & 0.472 \\
no warm-up & 93.8 & 7.1 & 0.656$\pm$0.049 & 0.547 & 0.497 \\
$\gamma=0$ (contextual bandit) & 94.4 & 6.4 & 0.645$\pm$0.070 & 0.576 & 0.510 \\
\bottomrule
\end{tabular}

\end{table*}

